\documentclass[letterpaper]{report}

\setlength\parindent{0pt}
\usepackage[utf8]{inputenc}
\usepackage[T1]{fontenc}
\usepackage{textcomp}
\usepackage{amsmath, amssymb}
\usepackage{enumitem}
\title{Weekly HW 3}
\author{Svadrut Kukunooru}
\date{\today}

\begin{document}
\begin{titlepage}
    \maketitle
\end{titlepage}

\subsection{Section 1.6, Problem 10}%
\label{sub:Section 1.6, Problem 10}
In the following problem, let $f : \mathbb{R}^7 \to \mathbb{R}^3$ for the matrix transofmration defined by $f(x) = Ax$, where 
\[
A = 
\begin{bmatrix} 
    1 & 2 \\
    0 & 1 \\
    1 & 1
\end{bmatrix} 
.\] 
Determine whether the matrix $w = \begin{bmatrix} 
1 \\ 1 \\ 1
\end{bmatrix}$ is in the range of $f$.
\[
\begin{bmatrix} 
    1 & 2 \\
    0 & 1 \\
    1 & 1
\end{bmatrix} 
\begin{bmatrix} 
x \\ y
\end{bmatrix} = 
\begin{bmatrix} 
1 \\ 1 \\ 1
\end{bmatrix} 
.\] 
\begin{equation}
    \begin{cases}
        x + 2y = 1 \\
        y = 1 \\
        x + y = 1
    \end{cases}
\end{equation}
Since this system of equations is impossible to solve, the matrix $w$ is NOT in the range of $f$. 

\subsection{Section 2.1, Problem 6}%
\label{sub:Section 2.1, Problem 6}
Find the reduced row echelon form of each of the given matrices. Record the row operations you perform, using the notation for elementary row operations. 

\begin{enumerate}[label=(\alph*)]
    \item \[\begin{bmatrix}
            -1 & 2 & -5 \\
            2 & -1 & 6 \\
            2 & -2 & 7
        \end{bmatrix}
    \] 
\end{enumerate}
Remember that reduced row echelon form is a matrix that has 3 characteristics: 
\begin{enumerate}
    \item The first nonzero entry in each row is 1 
    \item If a column has a leading 1, everything else is 0
    \item If a row has a leading 1, then every row has a leading 1 further to the left. 
\end{enumerate}
You can make RREFs using row operations. \\ \\ 
Multiply the first row by 2 and add it to the second and third row:
\[
\begin{bmatrix} 
    -1 & 2 & -5 \\
    0 & 3 & -4 \\
    2 & -2 & 7
\end{bmatrix} \to 
\begin{bmatrix} 
    1 & -2 & 5 \\
    0 & 3 & -4 \\
    0 & 2 & -3
\end{bmatrix} 
.\] 
Subtract the 2nd row from the 3rd row: 
\[
\begin{bmatrix} 
    1 & -2 & 5 \\
    0 & 3 & -4 \\
    0 & 2 & -3
\end{bmatrix} 
\to 
\begin{bmatrix} 
    1 & -2 & 5 \\
    0 & 1 & -1 \\
    0 & 2 & -3
\end{bmatrix} 
.\] 
Subtract two copies of 2nd row from 3rd row: 
\[
\begin{bmatrix} 
    1 & -2 & 5 \\
    0 & 1 & -1 \\
    0 & 2 & -3
\end{bmatrix} \to 
\begin{bmatrix} 
    1 & -2 & 5 \\
    0 & 1 & -1 \\
    0 & 0 & -1
\end{bmatrix}
.\] 
Add two copies of 2nd row to 1st row: 
\[
\begin{bmatrix} 
    1 & -2 & 5 \\
    0 & 1 & -1 \\
    0 & 0 & -1
\end{bmatrix}
\to 
\begin{bmatrix} 
    1 & 0 & 3 \\
    0 & 1 & -1 \\
    0 & 0 & -1  
\end{bmatrix} 
.\] 
Multiply the third row by -1: 
\[
\begin{bmatrix} 
    1 & 0 & 3 \\
    0 & 1 & -1 \\
    0 & 0 & -1
\end{bmatrix} \to 
\begin{bmatrix} 
    1 & 0 & 3 \\
    0 & 1 & -1 \\
    0 & 0 & 1
\end{bmatrix} 
.\] 
Add the third row to the second row: 
\[
\begin{bmatrix} 
    1 & 0 & 3 \\
    0 & 1 & -1 \\
    0 & 0 & 1
\end{bmatrix} \to 
\begin{bmatrix} 
    1 & 0 & 3 \\
    0 & 1 & 0 \\
    0 & 0 & 1
\end{bmatrix} 
.\] 
Add (negative) 3 copies of the third row to the first row: 
\[
\begin{bmatrix} 
    1 & 0 & 3 \\
    0 & 1 & 0 \\
    0 & 0 & 1
\end{bmatrix} \to 
\fbox{$
\begin{bmatrix} 
    1 & 0 & 0 \\
    0 & 1 & 0 \\
    0 & 0 & 1
\end{bmatrix} $}
.\] 
\begin{enumerate}[(label=\alph*)]
    \setcounter{enumi}(b)
    \item \[
    \begin{bmatrix} 
        1 & 1 & -1 \\
        3 & 4 & -1 \\
        5 & 6 & 3 \\
        -2 & -2 & 2
    \end{bmatrix} 
    .\] 
\end{enumerate}
I'm going to try another method for this problem, working from the bottom up rather than from the top down. 
\\ \\ 
q
\end{document}
